The following flow diagram illustrates the systematic literature
search and screening process following the PRISMA 2020 guidelines
\cite{Page2021}.  % replace or remove if not using BibTeX
 
\vspace{0.5cm}
 
\begin{figure}[h!]
\centering
\begin{tikzpicture}
 
% ────────────────────────────────────────────────────────────
%  All nodes use absolute coordinates so phase labels and
%  their content boxes are always vertically centred together.
%
%  Main-flow boxes (col x=0):  id1, id2, sc1–sc4, inc1
%  Phase labels   (col x=-5):  ph-id, ph-sc, ph-inc
%  Exclusion boxes(col x=8.2): exc1–exc4
% ────────────────────────────────────────────────────────────
 
% ── Row y-positions ─────────────────────────────────────────
%   id1  : y =  0
%   id2  : y = -2.0
%   sc1  : y = -4.2
%   sc2  : y = -6.0
%   sc3  : y = -7.8
%   sc4  : y = -9.6
%   inc1 : y = -11.8
 
% ═══════════════════════════════════════════════════════════
%  IDENTIFICATION
% ═══════════════════════════════════════════════════════════
 
% Phase label – vertically centred over id1 and id2
% Centre of id1=0, centre of id2=-2.0  →  mid = -1.0
\node[phase, minimum height=3.6cm] (ph-id) at (-5, -1.0)
  {Identi-\\fication};
 
\node[box] (id1) at (0, 0)
  {Records identified from:\\
   \textbf{Databases} (\textit{n}\,=\,\textlangle\textit{NUMBER}\textrangle)\\[1pt]
   \footnotesize e.g.\ PubMed, Scopus, Web of Science, \ldots};
 
\node[box] (id2) at (0, -2.0)
  {Records identified from other sources:\\
   \textbf{Registers / grey literature} (\textit{n}\,=\,\textlangle\textit{NUMBER}\textrangle)\\[1pt]
   \footnotesize e.g.\ ClinicalTrials.gov, hand-searching, \ldots};
 
\draw[arr] (id1.south) -- (id2.north);
 
% ═══════════════════════════════════════════════════════════
%  SCREENING
% ═══════════════════════════════════════════════════════════
 
% Phase label – centred over sc1(−4.2) through sc4(−9.6)  →  mid = −6.9
\node[phase, minimum height=6.8cm] (ph-sc) at (-5, -6.9)
  {Screening};
 
\node[box] (sc1) at (0, -4.2)
  {Records after duplicates removed\\
   (\textit{n}\,=\,\textlangle\textit{NUMBER}\textrangle)};
 
\node[excl, right=1.7cm of sc1] (exc1)
  {Records removed before screening:\\
   Duplicates removed (\textit{n}\,=\,\textlangle\textit{NUMBER}\textrangle)\\
   Flagged by automation tools (\textit{n}\,=\,\textlangle\textit{NUMBER}\textrangle)\\
   Other reasons (\textit{n}\,=\,\textlangle\textit{NUMBER}\textrangle)};
 
\draw[arr]  (id2.south) -- (sc1.north);
\draw[arr2] (sc1.east)  -- (exc1.west);
 
\node[box] (sc2) at (0, -6.0)
  {Records screened (title \& abstract)\\
   (\textit{n}\,=\,\textlangle\textit{NUMBER}\textrangle)};
 
\node[excl, right=1.7cm of sc2] (exc2)
  {Records excluded\\
   (\textit{n}\,=\,\textlangle\textit{NUMBER}\textrangle)};
 
\draw[arr]  (sc1.south) -- (sc2.north);
\draw[arr2] (sc2.east)  -- (exc2.west);
 
\node[box] (sc3) at (0, -7.8)
  {Reports sought for retrieval (full text)\\
   (\textit{n}\,=\,\textlangle\textit{NUMBER}\textrangle)};
 
\node[excl, right=1.7cm of sc3] (exc3)
  {Reports not retrieved\\
   (\textit{n}\,=\,\textlangle\textit{NUMBER}\textrangle)};
 
\draw[arr]  (sc2.south) -- (sc3.north);
\draw[arr2] (sc3.east)  -- (exc3.west);
 
\node[box] (sc4) at (0, -9.6)
  {Reports assessed for eligibility\\
   (\textit{n}\,=\,\textlangle\textit{NUMBER}\textrangle)};
 
\node[excl, right=1.7cm of sc4] (exc4)
  {Reports excluded:\\[2pt]
   \quad \textlangle\textit{Reason 1}\textrangle\ (\textit{n}\,=\,\textlangle\textit{N}\textrangle)\\
   \quad \textlangle\textit{Reason 2}\textrangle\ (\textit{n}\,=\,\textlangle\textit{N}\textrangle)\\
   \quad \textlangle\textit{Reason 3}\textrangle\ (\textit{n}\,=\,\textlangle\textit{N}\textrangle)\\
   \quad \ldots\\
   \quad Total excluded (\textit{n}\,=\,\textlangle\textit{NUMBER}\textrangle)};
 
\draw[arr]  (sc3.south) -- (sc4.north);
\draw[arr2] (sc4.east)  -- (exc4.west);
 
% ═══════════════════════════════════════════════════════════
%  INCLUDED
% ═══════════════════════════════════════════════════════════
 
% Phase label – centred on inc1 at y=−11.8
\node[phase, minimum height=2.0cm] (ph-inc) at (-5, -11.8)
  {Included};
 
\node[incl] (inc1) at (0, -11.8)
  {Studies included in review
   (\textit{n}\,=\,\textlangle\textit{NUMBER}\textrangle)\\[2pt]
   \footnotesize Reports of included studies
   (\textit{n}\,=\,\textlangle\textit{NUMBER}\textrangle)};
 
\draw[arr] (sc4.south) -- (inc1.north);
 
% ── Optional meta-analysis box – uncomment if needed ────────
%\node[incl] (inc2) at (0, -13.6)
%  {Studies included in meta-analysis\\
%   (\textit{n}\,=\,\textlangle\textit{NUMBER}\textrangle)};
%\draw[arr] (inc1.south) -- (inc2.north);
%\node[phase, minimum height=2.0cm] at (-5,-13.6) {Included};
 
\end{tikzpicture}
 
\caption{PRISMA 2020 flow diagram of the literature search and
         selection process.}
\label{fig:prisma}
\end{figure}
 
{\footnotesize
\textbf{Note.} Adapted from Page et al.\ (2021). \textit{The PRISMA
2020 statement: an updated guideline for reporting systematic
reviews.} \textit{BMJ}, 372, n71.
}